\section{早期证据：Machine Translation、Parsing 与 Attention Pattern}

机制合理并不等于模型有效。原始 Transformer 首先在 WMT machine translation 上展示质量与训练效率，再迁移到 constituency parsing；attention visualization 则提供局部可解释线索。本章把三类证据分开：benchmark improvement、cross-task transfer 和 qualitative pattern。

\subsection{Machine Translation：质量与训练效率同时改善}

原始论文在 WMT14 English--German 与 English--French 上比较 BLEU、训练成本和模型规模。读表时应先确认 single model/ensemble 与训练预算，再比较质量；Transformer 的关键不是只多零点几 BLEU，而是在更易并行的架构下达到或超过当时系统。

\lecturefigure{slide-24-sota-machine-translation.jpg}{原始 Transformer 的 machine-translation 结果}{官方录像恢复幻灯片 V024；00:28:54}

\begin{knowledgebox}{读图：先对齐模型类型和训练成本}
表格同时包含 recurrent、convolutional、single-model 与 ensemble baselines。先比较同类设置，再看训练 FLOPs/时间；关键证据是 Transformer 在较低训练成本下获得强 BLEU。BLEU 仍是 corpus-level overlap metric，不能覆盖所有语义、鲁棒性或真实部署质量。
\end{knowledgebox}

\subsection{Generalization to Parsing：架构不只适用于翻译}

Constituency parsing 的输入输出结构与 translation 不同。将同一 architecture 应用于 parsing，测试的是计算 motif 是否能跨任务复用，而不是只记住翻译特征。结果支持 self-attention + FFN stack 是一般 sequence transduction backbone。

\lecturefigure{slide-25-generalization-to-parsing.jpg}{Transformer 从 machine translation 泛化到 constituency parsing}{官方录像恢复幻灯片 V025；00:29:45}

\begin{knowledgebox}{读图：跨任务结果验证的是 architecture reuse}
先看 WSJ 等 parsing benchmark 与 baselines，再区分“同一 block 可训练”与“无需 task-specific data”。Parsing 仍有自己的输出 representation、数据与 objective；图支持 backbone consolidation，却不等于所有任务只需零样本 prompt。
\end{knowledgebox}

\subsection{Attention Map：可读 Pattern 不是完整解释}

部分 head 会展示长距离依赖、句法或复制模式，因此成为研究者理解模型的窗口。但 attention weight 只描述当前 layer/head 的路由系数，不包含 value content、residual contribution、后续 nonlinear transformation，也不能直接证明因果重要性。

\lecturefigure{slide-26-interpretable-attention-patterns.jpg}{Attention map 中可见的长距离与句法 pattern}{官方录像恢复幻灯片 V026；00:30:48}

\begin{warningbox}{“看起来像语法”不等于模型因果上依赖该 Head}
读图先确认横纵 token、颜色尺度和 head/layer，再观察高权重是否稳定跨样本。要证明机制，还需 ablation、activation patching、counterfactual intervention 或 value-path analysis；单张 heatmap 只能提供 hypothesis。
\end{warningbox}

\teachervoice{讲者对 attention interpretability 保持克制：有些 pattern 很漂亮，也确实像 long-distance dependency；但 architecture 的能力可能来自信息路由与巨大 FFN 的共同作用。不要把一张图升级为“模型思考过程”。}

\subsection{本章小结}

Early evidence 分成三层：translation 表格证明质量/效率，parsing 证明 backbone 可迁移，attention map 提供可研究的内部 pattern。三层都重要，却支持不同强度的结论；高质量讲义必须保留这种 evidence discipline。

